\section{Integral homology and geometric Morse counts}\label{sec:morse}

We first establish the topological comparison used in the construction.
Integral homology gives a lower bound for every Morse function that is exactly
quadratic outside a compact set.  On the simply connected manifolds of high
dimension used below, this bound is attained by an ordinary Morse function.
Both statements concern actual smooth functions.

For a closed smooth manifold $W$, write
\[
 r_i(W)=\rank H_i(W;\Z),\qquad
 t_i(W)=\text{the minimum number of generators of }\Tor H_i(W;\Z),
\]
and define
\begin{equation}\label{eq:lambda-definition}
 b(W)=\sum_i r_i(W),\qquad
 \lambda(W)=b(W)+2\sum_i t_i(W).
\end{equation}
All homological degrees in this section are integers.  We take $r_i(W)$ and
$t_i(W)$ to be zero outside the homological range of $W$.

Another convention allows a uniformly bounded differential perturbation
of a quadratic form at infinity.  It gives the same geometric minimum.

\begin{lemma}\label{lem:infinity-conventions}
Fix a Riemannian metric on a closed smooth manifold $W$ and the Euclidean
metric on $\R^k$, where $k>0$.  Suppose $Q$ is a nondegenerate quadratic
form on $\R^k$ and $F\colon W\times\R^k\to\R$ is a smooth function
with nondegenerate critical points such that
\[
 \sup_{W\times\R^k}\|dF-dQ\|<\infty
\]
in the product metric.  Then there is a smooth function $\widetilde F$
equal to $F$ near all its critical points, with no additional critical
points, and equal to $Q$ outside a compact set.  Consequently allowing
these bounded differential perturbations in the definition of the stable
Morse number does not change its value.
\end{lemma}

\begin{proof}
Write $h=F-Q$ and choose $K$ with $\|dh\|\leq K$.
Compactness of $W$ and integration along radial segments give
\[
 |h(x,z)|\leq K_0+K|z|,
 \qquad K_0=\max_{x\in W}|h(x,0)|.
\]
Nondegeneracy of $Q$ gives a constant $c>0$ with
$\|dQ(z)\|\geq c|z|$, so $F$ has no critical point for $|z|>K/c$.
Choose a smooth cutoff $\chi_R$ equal to one for $|z|\leq R$ and zero
for $|z|\geq2R$, with $0\leq\chi_R\leq1$ and
$\|d\chi_R\|\leq C/R$, where $C$ is independent of $R$.
For $R\geq1$, the differential of $\chi_Rh$ on the intervening annulus
has norm at most
\[
 K+C(K_0+2K).
\]
This bound is independent of $R$, whereas $\|dQ\|\geq cR$ there.
For sufficiently large $R$, the function
$\widetilde F=Q+\chi_Rh$ therefore has no critical point in the annulus.
It agrees with $F$ on a ball containing all the critical points and
with $Q$ outside the larger ball.  Conversely, equality with $Q$ outside
a compact set implies the displayed boundedness condition.  These two
constructions preserve the number of critical points; for $k=0$ both
conventions allow the same ordinary Morse functions.
\end{proof}

\begin{proposition}\label{prop:stable-lower}
For every connected closed smooth manifold $W$,
\[
 \SM(W)\geq \lambda(W).
\]
\end{proposition}

\begin{proof}
Let $F\colon W\times\R^k\to\R$ be an admissible function in the definition
of $\SM(W)$.  We construct a finite free integral chain complex with one
generator for each critical point of $F$ and with homology equal to that of
$W$, shifted by a fixed integer.  When $k=0$, the ordinary Morse complex
has these properties.  Suppose $k>0$.

After a linear change of the auxiliary coordinates, the quadratic form is
\[
 Q(u,v)=-|u|^2+|v|^2,
 \qquad (u,v)\in\R^d\times\R^{k-d}.
\]
The equality $F=Q$ still holds outside a sufficiently large Euclidean ball.
Choose a radius $R_0$ such that this equality holds for $|(u,v)|\geq R_0$
and all critical points lie in $W\times B_{R_0}$.  Choose $A>0$ satisfying
\[
 A>\max_{W\times\overline B_{R_0}}\{|F|,|Q|\},
\]
and then choose $T>R_0$ with $T^2>A$.  At the two levels $-A$ and $A$,
the sublevel sets of $F$ and $Q$ agree, including inside the core
$W\times\overline B_{R_0}$: at $-A$ the core is excluded, and at $A$ it
is included.  Thus, with
\[
 K=W\times\overline B_T,\qquad
 K_-=K\cap\{F\leq-A\},\qquad
 K_+=K\cap\{F\leq A\},
\]
we have $K_\pm=K\cap\{Q\leq\pm A\}$.

This truncation introduces no boundary critical point in the level band
$[-A,A]$.  Indeed, on the side boundary $W\times\partial B_T$ the function
is $Q$, and its restricted differential can vanish only when $u=0$ or
$v=0$.  Its possible critical values there are therefore $T^2$ and $-T^2$.
In the band, a gradient-like vector field can consequently be chosen
tangent to the side boundary.  The usual Morse attachment argument on
this compact manifold with corners gives $K_+$ the relative homotopy type,
over $K_-$, of a space obtained by attaching one cell of the appropriate
Morse index for each critical point of $F$;
see \cite[Section~3]{Milnor1963} for the attachment theorem.
The tangent field gives the same deformation argument at the side
boundary as in the interior.  If necessary, separate the critical values
by a small perturbation supported near the critical points.  The endpoint
sublevels are compact manifolds with corners, or empty, and have finite
CW type.  Taking a finite CW model for $K_-$ and using cellular approximation for the
attachments yields a finite free relative cellular complex $C_*$ with
\begin{equation}\label{eq:critical-chain-rank}
 \sum_i\rank C_i=\#\operatorname{Crit}(F),\qquad
 H_*(C_*)\cong H_*(K_+,K_-;\Z).
\end{equation}

The endpoint pair has a simple geometric description.  The homotopy
$(u,v)\mapsto(u,(1-s)v)$, $0\leq s\leq1$, decreases both $Q$ and the
Euclidean norm.  It therefore preserves both members of the pair and
deformation retracts $(K_+,K_-)$ onto
\[
 W\times\bigl(D_T^d,
       \{u\in D_T^d:|u|\geq\sqrt A\}\bigr).
\]
For $d=0$ the second factor is the pair $(\{0\},\varnothing)$.
The relative integral homology of the displayed disk--annulus pair is
$\Z$ in degree $d$ and zero otherwise.  The relative K\"unneth formula
therefore gives
\begin{equation}\label{eq:quadratic-pair-homology}
 H_i(C_*)\cong H_{i-d}(W;\Z).
\end{equation}
This also covers positive or negative definite quadratic forms.

It remains to count the ranks in $C_*$.  For any finite free integral
chain complex, put $Z_i=\ker\partial_i$ and
$B_i=\operatorname{im}\partial_{i+1}$.  Since subgroups of free abelian
groups are free,
\begin{equation}\label{eq:chain-rank}
 \rank C_i=\rank H_i(C_*)+\rank B_i+\rank B_{i-1}.
\end{equation}
Smith normal form for $B_i\subset Z_i$ shows that $\rank B_i$ is at
least the minimum number of generators of $\Tor H_i(C_*)$.
Summing \eqref{eq:chain-rank} and applying
\eqref{eq:critical-chain-rank}--\eqref{eq:quadratic-pair-homology} gives
$\#\operatorname{Crit}(F)\geq\lambda(W)$.  Taking the minimum over
admissible $F$ proves the proposition.
\end{proof}

For the upper bound we use the following form of Smale's Morse theorem.
We state the dimension range for which the handle proof in
Appendix~\ref{app:handles} is given.

\begin{theorem}[Smale]\label{thm:smale}
Let $W$ be a simply connected closed smooth manifold of dimension $D\geq9$.
There is an ordinary Morse function on $W$ with
\[
 r_i(W)+t_i(W)+t_{i-1}(W)
\]
critical points of index $i$ for every $i$.  In particular,
\begin{equation}\label{eq:simply-connected-morse}
 \SM(W)=\Morse(W)=\lambda(W).
\end{equation}
\end{theorem}

\begin{proof}
The existence statement is the torsion-inclusive form of Smale's theorem;
see \cite[Theorem~2.4 and the following paragraph]{Smale1963}.
Appendix~\ref{app:handles} supplies a handle proof in the stated range,
including the cancellation of excess handles in degrees $2$ and $3$.
Summing the prescribed index counts gives an ordinary Morse function with
$\lambda(W)$ critical points.  Proposition~\ref{prop:stable-lower} and
the inclusion of ordinary Morse functions among the admissible functions
then give
\[
 \lambda(W)\leq\SM(W)\leq\Morse(W)\leq\lambda(W).
\]
\end{proof}

\begin{remark}\label{rem:torsion-generators}
The number $t_i(W)$ counts cyclic invariant factors, so it can combine
torsion at different primes.  For example, for distinct primes $p,q$,
the minimum number of generators of
$(\Z/p)^a\oplus(\Z/q)^b$ is $\max(a,b)$.
Reduction modulo $p$ and modulo $q$ gives the lower bound; pairing
$\min(a,b)$ summands by the Chinese remainder theorem gives the upper
bound.  This dependence on the degreewise distribution of primes is the
feature of $\lambda$ used below.
\end{remark}
